\section{Problem Formulation}
\label{sec:problem_formulation}

This thesis formulates network intrusion detection as hierarchical binary
classification over two levels of network structure. At the first level, an
individual flow is represented by its ordered packet observations and a
separate vector of flow-level statistics. At the second level, the target flow
is interpreted in relation to preceding flows selected by temporal and endpoint
metadata. This formulation follows the observation that packet timing can carry
discriminative information within a connection \cite{Han2023GTID}, while
large-scale attacks can exhibit dependencies across multiple network events
\cite{LiuPatras2022NetSentry,Manocchio2024FlowTransformer}. The objective is to
predict a label for every target flow without inspecting plaintext payload and
without using information from future flows.

\subsection{Scope and Assumptions}
\label{subsec:pf_scope}

The detector operates on traffic reconstructed from PCAP files. Packet- and
flow-level records are produced by a custom Suricata output plugin, after which
the learning pipeline performs flow-level binary classification. Label $0$
denotes benign traffic and label $1$ denotes malicious traffic. The formulation
does not assume access to application payload semantics; packet-header fields,
directional metadata, packet timing, and flow statistics form the model input.
Suricata alerts may be used to derive supervision, but alert indicators, rule
identifiers, and other direct label proxies are excluded from the feature space.
This separation is necessary to prevent the learned detector from reproducing
the labeling mechanism rather than learning traffic behavior.

The following assumptions define the scope of the problem:

\begin{itemize}
    \item each packet is assigned to one bidirectional flow through a stable
    flow identifier;
    \item packets inside a flow are ordered by capture timestamp;
    \item flows are ordered by their start timestamp before Stage2 context is
    constructed;
    \item only the current flow and eligible preceding flows may be used to
    classify a target flow;
    \item train-fitted preprocessing parameters are reused for validation and
    test data; and
    \item the output is a flow-level binary decision rather than packet-level
    localization or attack-family attribution.
\end{itemize}

These assumptions reflect a passive NIDS setting. The model detects suspicious
flows but does not itself block traffic or replace rule-based inspection.
Instead, it complements systems such as Suricata by learning behavioral
patterns that may not be represented by a fixed signature. This distinction is
important because operational network traffic is an open-world domain and
laboratory performance does not automatically imply deployment performance
\cite{SommerPaxson2010}.

\subsection{Input Space and Notation}
\label{subsec:pf_notation}

Let the chronologically ordered flow dataset be

\begin{equation}
\mathcal{D}
=
\left\{
\left(\mathcal{P}_i,\mathbf{s}_i,\mathbf{m}_i,y_i\right)
\right\}_{i=1}^{N},
\label{eq:dataset_definition}
\end{equation}

where $N$ is the number of flows and $y_i\in\{0,1\}$ is the label of flow
$i$. The four components associated with the flow have different roles.

First, the packet sequence is

\begin{equation}
\mathcal{P}_i
=
\left(
(\mathbf{x}_{i,1}^{p},t_{i,1}),
\ldots,
(\mathbf{x}_{i,T_i}^{p},t_{i,T_i})
\right),
\label{eq:packet_sequence}
\end{equation}

where $T_i$ is the number of captured packets, $\mathbf{x}_{i,k}^{p}\in
\mathbf{R}^{d_p}$ is the payload-independent packet feature vector, and
$t_{i,k}$ is the packet timestamp. The packet feature vector can contain
numerical, categorical, and binary header-derived attributes after
preprocessing.

Second, $\mathbf{s}_i\in\mathbf{R}^{d_f}$ is a flow-statistical vector. It
summarizes properties such as duration, packet and byte counts, rates, and
inter-arrival-time statistics. Crucially, $\mathbf{s}_i$ is not repeated and
concatenated to every packet token in the principal formulation. It is encoded
through a separate branch and combined with the packet-sequence representation
by late fusion.

Third, the metadata tuple

\begin{equation}
\mathbf{m}_i=(u_i,v_i,t_i^{\mathrm{start}},q_i)
\label{eq:flow_metadata}
\end{equation}

contains the source identifier $u_i$, destination identifier $v_i$, flow start
time $t_i^{\mathrm{start}}$, and split identifier $q_i\in
\{\mathrm{train},\mathrm{validation},\mathrm{test}\}$. These fields determine
which flows may enter the Stage2 context. Unless explicitly included in a
controlled ablation, they are indexing metadata rather than direct classifier
features; this reduces the risk that the classifier simply memorizes host
identities.

Because flows have variable packet counts, a sequence selection operator
$\mathcal{S}_{L}$ maps $\mathcal{P}_i$ to at most $L$ packets:

\begin{equation}
\widetilde{\mathcal{P}}_i
=
\mathcal{S}_{L}(\mathcal{P}_i),
\qquad
\widetilde{T}_i=\min(T_i,L).
\label{eq:sequence_selection}
\end{equation}

The operator may retain the first packets, combine the head and tail, or sample
packets according to a specified strategy. Shorter sequences are padded, and a
mask $\mathbf{M}_i^{p}\in\{0,1\}^{L}$ distinguishes observed packets from
padding. The selection strategy and $L$ are experimental variables rather than
part of the class label.

\subsection{Temporal Representation Within a Flow}
\label{subsec:pf_time}

Ordinal packet position and physical network time represent different
information. Position identifies whether one packet precedes another, whereas
elapsed time distinguishes a rapid burst from a sequence containing long idle
gaps. For consecutive packets, define the non-negative inter-arrival time

\begin{equation}
\Delta t_{i,k}
=
\begin{cases}
0, & k=1,\\
\max\left(0,t_{i,k}-t_{i,k-1}\right), & k>1,
\end{cases}
\label{eq:packet_iat}
\end{equation}

and the cumulative elapsed time

\begin{equation}
c_{i,k}=\sum_{r=1}^{k}\Delta t_{i,r}.
\label{eq:cumulative_time}
\end{equation}

Since network timing is typically heavy-tailed, the stable time variables are

\begin{equation}
\rho_{i,k}=\log(1+\Delta t_{i,k}),
\qquad
\kappa_{i,k}=\log(1+c_{i,k}+\epsilon),
\label{eq:log_time}
\end{equation}

where $\epsilon>0$ ensures numerical stability. The local variable
$\rho_{i,k}$ represents burst-level timing, while $\kappa_{i,k}$ preserves
longer-range elapsed time. Explicit use of inter-packet intervals is motivated
by time-aware traffic models such as GTID \cite{Han2023GTID}; it is distinct
from standard positional encoding, which represents order but not the physical
distance between observations \cite{Vaswani2017}.

\subsection{Stage1: Intra-Flow Representation Learning}
\label{subsec:pf_stage1}

Stage1 learns a fixed-dimensional representation from the variable-length
packet sequence and the flow-statistical vector. Each selected packet is first
projected into a latent space:

\begin{equation}
\mathbf{e}_{i,k}=\mathbf{W}_{p}\mathbf{x}_{i,k}^{p}+\mathbf{b}_{p},
\qquad
\mathbf{e}_{i,k}\in\mathbf{R}^{d}.
\label{eq:packet_projection}
\end{equation}

Position and time are then incorporated through a configurable encoding
function:

\begin{equation}
\widetilde{\mathbf{e}}_{i,k}
=
\Phi_{\mathrm{pt}}
\left(
\mathbf{e}_{i,k},k,\rho_{i,k},\kappa_{i,k}
\right).
\label{eq:position_time_encoding}
\end{equation}

The function $\Phi_{\mathrm{pt}}$ may add, concatenate, or gate independent
position and time encodings. This abstraction allows the empirical comparison
of no encoding, position-only encoding, time-only encoding, and combined
encoding while leaving the packet features unchanged.

The encoded sequence is processed by an intra-flow encoder
$E_{\theta_1}$, implemented by the principal model as a Transformer:

\begin{equation}
\mathbf{H}_{i}^{p}
=
E_{\theta_1}
\left(
\widetilde{\mathbf{e}}_{i,1:L},
\mathbf{M}_{i}^{p}
\right),
\qquad
\mathbf{H}_{i}^{p}\in\mathbf{R}^{L\times d}.
\label{eq:stage1_transformer}
\end{equation}

Self-attention permits direct interaction among packet tokens and has been
shown to be effective for sequential network data
\cite{Vaswani2017,Manocchio2024FlowTransformer}. Padding positions are excluded
through the mask. A masked attention-pooling operator produces a packet-derived
flow representation:

\begin{align}
a_{i,k}
&=
\mathbf{w}_{a}^{\top}
\tanh(\mathbf{W}_{a}\mathbf{h}_{i,k}^{p}+\mathbf{b}_{a}),
\label{eq:attention_score}\\
\alpha_{i,k}
&=
\frac{
M_{i,k}^{p}\exp(a_{i,k})
}{
\sum_{r=1}^{L}M_{i,r}^{p}\exp(a_{i,r})
},
\label{eq:attention_weight}\\
\mathbf{z}_{i}^{p}
&=
\sum_{k=1}^{L}\alpha_{i,k}\mathbf{h}_{i,k}^{p}.
\label{eq:packet_embedding}
\end{align}

In parallel, the flow statistics are mapped into the same latent dimension:

\begin{equation}
\mathbf{z}_{i}^{f}=G_{\theta_f}(\mathbf{s}_i),
\qquad
\mathbf{z}_{i}^{f}\in\mathbf{R}^{d}.
\label{eq:flow_embedding}
\end{equation}

The final Stage1 embedding is obtained by late fusion:

\begin{equation}
\mathbf{z}_{i}^{(1)}
=
\Psi_{\theta_{\mathrm{fus}}}
\left(
\mathbf{z}_{i}^{p},\mathbf{z}_{i}^{f}
\right).
\label{eq:late_fusion}
\end{equation}

For example, a gated fusion can be written as

\begin{align}
\mathbf{g}_{i}
&=
\sigma\left(
\mathbf{W}_{g}
[\mathbf{z}_{i}^{p};\mathbf{z}_{i}^{f}]
+\mathbf{b}_{g}
\right),
\label{eq:fusion_gate}\\
\mathbf{z}_{i}^{(1)}
&=
\mathbf{g}_{i}\odot\mathbf{z}_{i}^{p}
+
(1-\mathbf{g}_{i})\odot\mathbf{z}_{i}^{f},
\label{eq:gated_fusion}
\end{align}

where $[\cdot;\cdot]$ denotes concatenation and $\odot$ denotes element-wise
multiplication. Concatenative and additive fusion are alternative
instantiations of $\Psi$. An auxiliary Stage1 classifier produces

\begin{equation}
\mathbf{p}_{i}^{(1)}
=
\operatorname{softmax}
\left(g_{1}(\mathbf{z}_{i}^{(1)})\right).
\label{eq:stage1_probability}
\end{equation}

The Stage1 classification objective makes the exported embedding directly
informative for flow-level detection. It also permits an independent test of
RQ1 before any cross-flow information is introduced.

\subsection{Stage2 Context Construction}
\label{subsec:pf_context}

Stage2 is defined over a causal sequence of Stage1 embeddings. Let
$\mathcal{H}_i=\{j: j<i\}$ denote the set of flows observed before target flow
$i$ in global chronological order. Four relation indicators define the context
construction methods:

\begin{align}
R_{\mathrm{time}}(i,j)
&=1,
\label{eq:relation_time}\\
R_{\mathrm{src}}(i,j)
&=\mathbf{1}[u_i=u_j],
\label{eq:relation_source}\\
R_{\mathrm{dst}}(i,j)
&=\mathbf{1}[v_i=v_j],
\label{eq:relation_destination}\\
R_{\mathrm{end}}(i,j)
&=\mathbf{1}
\left[
\{u_i,v_i\}\cap\{u_j,v_j\}\neq\varnothing
\right].
\label{eq:relation_endpoint}
\end{align}

Here $\mathbf{1}[\cdot]$ is an indicator function. For method
$r\in\{\mathrm{time},\mathrm{src},\mathrm{dst},\mathrm{end}\}$, the eligible
history is

\begin{equation}
\mathcal{A}_{i}^{r}
=
\left\{
j\in\mathcal{H}_i:
R_{r}(i,j)=1
\;\land\;
\Pi(i,j)=1
\right\},
\label{eq:eligible_context}
\end{equation}

where $\Pi(i,j)$ is an evaluation-policy constraint. Under a
\emph{split-isolated} policy, $\Pi(i,j)=1$ only when $q_i=q_j$. Under an
\emph{online} policy, any earlier flow may be used without accessing its label.
Under a \emph{train-only-for-evaluation} policy, validation and test targets may
retrieve only eligible training flows. Making this policy explicit is necessary
because context construction can otherwise introduce cross-split or future
information, a known source of optimistic evaluation in machine-learning
security research \cite{Arp2022DosDonts}.

Let $\operatorname{Recent}_{W-1}(\cdot)$ return the most recent $W-1$ eligible
indices in chronological order. With the target included as the final token,
the context index sequence is

\begin{equation}
\mathcal{C}_{i}^{r,W}
=
\operatorname{Recent}_{W-1}(\mathcal{A}_{i}^{r})
\mathbin{\Vert}
(i),
\label{eq:context_window}
\end{equation}

where $\Vert$ denotes ordered sequence concatenation. This formulation never
uses $j>i$ and therefore prevents future-flow leakage. Source-, destination-,
and endpoint-conditioned histories represent relational context, whereas
time-only history provides a control that uses temporal proximity without host
matching. Related work has shown the value of both temporal flow sequences and
relational traffic structures
\cite{Deng2023FlowTopology,MartinReizabal2026FlowWindow}.

\subsection{Stage2 Context-Aware Detection}
\label{subsec:pf_stage2}

For context sequence $\mathcal{C}_{i}^{r,W}=(j_1,\ldots,j_m,i)$, define

\begin{equation}
\mathbf{Z}_{i}^{r,W}
=
\left[
\mathbf{z}_{j_1}^{(1)},
\ldots,
\mathbf{z}_{j_m}^{(1)},
\mathbf{z}_{i}^{(1)}
\right].
\label{eq:stage2_input}
\end{equation}

The sequence is padded to length $W$ when necessary, and
$\mathbf{M}_{i}^{c}\in\{0,1\}^{W}$ identifies valid flow tokens. A context
encoder $C_{\theta_2}$ maps the sequence to context-aware states:

\begin{equation}
\mathbf{H}_{i}^{c}
=
C_{\theta_2}
\left(
\mathbf{Z}_{i}^{r,W},
\mathbf{M}_{i}^{c}
\right).
\label{eq:stage2_encoder}
\end{equation}

The principal Stage2 model implements $C_{\theta_2}$ with Transformer
self-attention, while LSTM, GRU, and CNN--LSTM encoders provide architectural
comparisons under the same context indices. A no-context classifier uses only
$\mathbf{z}_{i}^{(1)}$ and tests whether improvements arise from historical
evidence rather than from another classification head.

Let $\operatorname{TargetPool}(\cdot)$ select or aggregate the target-aware
representation. Since the target is the last valid token,

\begin{equation}
\mathbf{z}_{i}^{(2)}
=
\operatorname{TargetPool}
\left(
\mathbf{H}_{i}^{c},\mathbf{M}_{i}^{c}
\right),
\label{eq:stage2_embedding}
\end{equation}

and the malicious-flow probability is

\begin{equation}
p_i
=
\left[
\operatorname{softmax}
\left(g_{2}(\mathbf{z}_{i}^{(2)})\right)
\right]_{1}.
\label{eq:final_probability}
\end{equation}

Thus, the complete hierarchical detector can be written as

\begin{equation}
p_i
=
f_{\theta}
\left(
\mathcal{P}_i,
\mathbf{s}_i,
\{(\mathcal{P}_j,\mathbf{s}_j):j\in\mathcal{C}_{i}^{r,W}\setminus\{i\}\}
\right),
\label{eq:hierarchical_detector}
\end{equation}

with $\theta=(\theta_1,\theta_f,\theta_{\mathrm{fus}},\theta_2)$. The hierarchy
separates two questions: whether packet timing improves an individual flow
representation, and whether preceding related flows improve the classification
of the target.

\subsection{Learning Objective Under Class Imbalance}
\label{subsec:pf_objective}

Stage1 and Stage2 are trained sequentially in the implemented pipeline. Stage1
is first optimized on flow labels, after which its fused embeddings are
exported. Stage2 is then trained on the exported embeddings and the context
indices. Let $\mathcal{I}_{\mathrm{tr}}$ denote training-flow indices. The two
optimization problems are

\begin{align}
\theta_{1}^{*}
&=
\arg\min_{\theta_1,\theta_f,\theta_{\mathrm{fus}}}
\frac{1}{|\mathcal{I}_{\mathrm{tr}}|}
\sum_{i\in\mathcal{I}_{\mathrm{tr}}}
\ell_{1}
\left(y_i,\mathbf{p}_{i}^{(1)}\right)
+\lambda_1\|\theta_{1}\|_{2}^{2},
\label{eq:stage1_objective}\\
\theta_{2}^{*}
&=
\arg\min_{\theta_2}
\frac{1}{|\mathcal{I}_{\mathrm{tr}}|}
\sum_{i\in\mathcal{I}_{\mathrm{tr}}}
\ell_{2}(y_i,p_i)
+\lambda_2\|\theta_{2}\|_{2}^{2}.
\label{eq:stage2_objective}
\end{align}

Because malicious flows are the minority class, $\ell_1$ and $\ell_2$ must not
allow benign examples to dominate optimization. A class-weighted
cross-entropy loss is

\begin{equation}
\mathcal{L}_{\mathrm{WCE}}
=
-\frac{1}{B}
\sum_{i=1}^{B}
\omega_{y_i}\log p_{i,y_i},
\label{eq:weighted_cross_entropy}
\end{equation}

where $B$ is the batch size and $\omega_{y_i}$ is a class weight. An alternative
is focal loss \cite{Lin2017FocalLoss}:

\begin{equation}
\mathcal{L}_{\mathrm{focal}}
=
-\frac{1}{B}
\sum_{i=1}^{B}
\alpha_{y_i}
\left(1-p_{i,y_i}\right)^{\gamma}
\log p_{i,y_i},
\label{eq:focal_loss}
\end{equation}

where $\gamma\geq0$ reduces the contribution of easy examples. Class-aware
sampling may additionally alter the training distribution, but it does not
change the validation or test distribution. The exact loss and sampling policy
must therefore be reported with each experiment rather than being treated as an
implicit property of the architecture.

\subsection{Decision Rule and Evaluation Objective}
\label{subsec:pf_evaluation}

The probability $p_i$ is converted to a binary decision using threshold
$\tau$:

\begin{equation}
\widehat{y}_i(\tau)=\mathbf{1}[p_i\geq\tau].
\label{eq:decision_rule}
\end{equation}

If threshold optimization is used, it is performed only on validation data:

\begin{equation}
\tau^{*}
=
\arg\max_{\tau\in\mathcal{T}}
Q_{\mathrm{val}}(\tau),
\label{eq:threshold_selection}
\end{equation}

where $Q_{\mathrm{val}}$ is a predeclared validation metric and
$\mathcal{T}\subset(0,1)$ is the search grid. The selected $\tau^{*}$ is then
fixed before test evaluation. Selecting it on test labels would leak test-set
information into the final result \cite{Arp2022DosDonts}.

For binary labels, the confusion matrix is written as

\begin{equation}
\mathbf{C}
=
\begin{bmatrix}
\mathrm{TN} & \mathrm{FP}\\
\mathrm{FN} & \mathrm{TP}
\end{bmatrix},
\label{eq:confusion_matrix}
\end{equation}

and the operational false positive rate is

\begin{equation}
\mathrm{FPR}
=
\frac{\mathrm{FP}}{\mathrm{FP}+\mathrm{TN}}.
\label{eq:fpr}
\end{equation}

For class $c\in\{0,1\}$, let $P_c$, $R_c$, and $F1_c$ denote precision,
recall, and F1. The principal macro metrics are

\begin{equation}
P_{\mathrm{macro}}=\frac{P_0+P_1}{2},
\qquad
R_{\mathrm{macro}}=\frac{R_0+R_1}{2},
\qquad
F1_{\mathrm{macro}}=\frac{F1_0+F1_1}{2}.
\label{eq:macro_metrics}
\end{equation}

Macro-F1 is the primary summary metric because it gives equal importance to
benign and malicious classes. Attack-class precision, recall, and F1 explain
the minority-class behavior directly. ROC-AUC and PR-AUC measure ranking
quality across thresholds, with PR-AUC receiving particular emphasis because
precision--recall analysis is more informative under strong class imbalance
\cite{SaitoRehmsmeier2015}. Accuracy and weighted metrics are reported as
secondary measures because they can be dominated by the benign class.

The desired detector therefore maximizes macro-F1 and PR-AUC while maintaining
a low FPR and acceptable computational cost. These criteria are complementary:
a model may increase attack recall by generating substantially more false
alarms, and a small gain in threshold-independent AUC may not translate into a
better operating point.

\subsection{Research Questions and Testable Hypotheses}
\label{subsec:pf_hypotheses}

The formalization leads to the following research questions and hypotheses.

\paragraph{RQ1: Intra-flow representation.}
Does explicit modeling of packet order and elapsed time produce a more
informative flow representation than flow-level statistics alone?

\begin{itemize}
    \item \textbf{H1a:} the combined position-and-time Stage1 model achieves
    higher macro-F1 and PR-AUC than a flow-statistics MLP;
    \item \textbf{H1b:} explicit time encoding improves over no encoding and
    position-only encoding; and
    \item \textbf{H1c:} late fusion of packet and flow representations improves
    over either representation alone.
\end{itemize}

\paragraph{RQ2: Inter-flow context.}
Does causal context from related preceding flows improve target-flow intrusion
detection compared with independent flow classification?

\begin{itemize}
    \item \textbf{H2a:} the context-aware Stage2 model improves macro-F1 and
    PR-AUC over the no-context classifier;
    \item \textbf{H2b:} host- or endpoint-conditioned context is more useful
    than time-only context when attack behavior is distributed across related
    communications; and
    \item \textbf{H2c:} context window size exhibits a finite optimum because
    short windows omit evidence while excessively long windows introduce
    unrelated traffic.
\end{itemize}

The hypotheses are evaluated through controlled ablations that keep the data
split and all unrelated settings fixed. Since models evaluated on different
test populations are not strictly comparable, conclusions are drawn only from
experiments sharing the same samples and labeling protocol.

\subsection{Problem Summary}
\label{subsec:pf_summary}

Given a payload-independent packet sequence $\mathcal{P}_i$, flow statistics
$\mathbf{s}_i$, and a causal relation-constrained history
$\mathcal{C}_{i}^{r,W}$, the problem is to learn a function that predicts the
binary label of flow $i$ while satisfying four requirements:

\begin{itemize}
    \item preserve packet order and physical elapsed time within each flow;
    \item combine packet dynamics with complementary flow statistics;
    \item use only eligible preceding flows for cross-flow reasoning; and
    \item remain reliable under class imbalance, low attack prevalence, and
    operational false-positive constraints.
\end{itemize}

The next chapter instantiates this formulation as a reproducible PCAP-to-model
pipeline, describing Suricata-based feature extraction, preprocessing, the
Stage1 time-aware Transformer, late fusion, Stage2 context construction, model
training, and evaluation.
